\documentclass[11pt]{article}

\usepackage[margin=1in]{geometry}
\usepackage{amsmath}
\usepackage{amssymb}
\usepackage{graphicx}
\usepackage[hidelinks]{hyperref}

\title{Where Does the Signal Die?\\
\large Localizing and Repairing the Temporal-Integration Bottleneck in
Deep Low-SNR Drone Detection}
\author{Roey Graif}
\date{May 2026}

\begin{document}

\maketitle

\begin{abstract}
\noindent
Deep networks detect drones reliably at moderate signal levels, but their
accuracy falls quickly when the target is faint, as it is for a small drone seen
from a distance by a thermal infrared camera. When such a detector fails, one
usually cannot tell whether the target was too weak for any method to find, or
whether the network had the evidence and threw it away. We separate these two
cases using synthetic data, for which the optimal Neyman--Pearson detector can be
computed exactly and used as a reference. We measure detectability at each stage
inside the network and compare it with this optimum. The stage where
detectability drops is the one we redesign, and we then test whether the lost
performance returns. The same setup shows how far integration over many frames
keeps helping, and what stops it: in a real thermal sensor a fixed noise pattern
that does not average out across frames sets a limit that depends on how fast the
target moves. Using the optimal detector as a reference is an idea borrowed from
medical imaging; applying it inside a deep detector, to find and repair where a
faint moving target is lost, appears to be new.
\end{abstract}

\section{Introduction}

Detecting small drones at long range increasingly relies on thermal infrared
cameras. At long range a drone covers only a few pixels and is often dimmer than
the background noise, the regime known as low signal-to-noise ratio (SNR).
Multi-frame deep networks, which combine evidence across a sequence of frames,
are the standard tool for this task, yet their probability of detection drops
sharply as the SNR falls.

When a deep detector misses a faint drone, there are two very different reasons
it might have failed. The target may carry too little information to be recovered
by any method, in which case the data sets the limit. Alternatively, the
information may be present in the frames while the network discards it somewhere
between its input and its decision, in which case the architecture sets the limit
and can be improved. On real footage these two cases cannot be told apart,
because separating them needs the optimal detector, and the optimal detector
cannot be computed when the signal and noise are not known exactly.

This proposal removes that obstacle by working in a synthetic infrared testbed
where the signal and the noise are known by construction. For a known signal in
Gaussian noise, the optimal detector in the Neyman--Pearson sense is the matched
filter, which becomes a whitening matched filter when the noise is correlated;
in both cases it can be computed \cite{kay,marcum}. It gives an exact upper bound
on detectability, against which a learned detector and each of its internal
representations can be measured. This use of the optimal detector as a reference
is standard in medical imaging, where the ``ideal observer'' tells whether image
processing has preserved a faint signal \cite{barrett}. Here the same idea is
applied one level deeper, inside the network.

\paragraph{Aim and research questions.} The thesis asks where, and why, a deep
low-SNR drone detector falls short of the optimum, and whether the shortfall can
be repaired:
\begin{enumerate}
\item[Q1.] Is there a stage in the network where detectability collapses relative
to the optimum, and if so, where is it?
\item[Q2.] Does redesigning that stage recover detectability toward the optimum,
by an amount consistent with theory?
\item[Q3.] How far does adding frames keep improving detection before
fixed-pattern sensor noise caps the gain, and how does this depend on SNR,
sequence length, and target speed?
\end{enumerate}

\paragraph{Contributions.} The planned contributions are:
\begin{enumerate}
\item[C1.] a method that plots detectability against network depth, with the
optimal detector as the ceiling;
\item[C2.] a map of where low-SNR signal is lost inside a deep
track-before-detect detector, as a function of SNR, sequence length, and motion;
\item[C3.] a redesign of the stage responsible for the loss, with a test of how
much detectability it recovers;
\item[C4.] a description of the limit that temporal integration reaches once
fixed-pattern noise is present.
\end{enumerate}

\section{Problem Formulation and Background}

\paragraph{Observation model.} Let $y_t \in \mathbb{R}^{H\times W}$ be frame $t$
of a sequence of $T$ frames. Under the target-present ($H_1$) and target-absent
($H_0$) hypotheses,
\begin{equation}
H_1:\; y_t = a_t\, g(\mathbf{p}_t;\sigma_{\mathrm{PSF}}) + n_t,
\qquad
H_0:\; y_t = n_t,
\qquad t = 1,\dots,T,
\label{eq:model}
\end{equation}
where $g(\cdot;\sigma_{\mathrm{PSF}})$ is a unit-peak Gaussian point-spread
function centered at the sub-pixel position $\mathbf{p}_t$ on the target
trajectory, $a_t$ is the per-frame amplitude, and $n_t$ is the sensor noise. We
measure target strength by the peak SNR,
$\mathrm{SNR}_{\mathrm{dB}} = 20\log_{10}(A/\sigma)$, with $A$ the peak target
amplitude and $\sigma$ the noise standard deviation.

\paragraph{The matched filter and the $\sqrt{T}$ rule.} Collect the frames and
the known signal template along the trajectory into vectors $y$ and $s$. For
white Gaussian noise $n\sim\mathcal{N}(0,\sigma^2 I)$ the optimal test is the
matched filter $\Lambda(y) = s^\top y$ \cite{kay}. Its deflection coefficient,
the squared separation between the two hypotheses divided by the variance, is
\begin{equation}
d^2 \;=\; \frac{\lVert s\rVert^2}{\sigma^2}
\;=\; \frac{1}{\sigma^2}\sum_{t=1}^{T} a_t^2\,\lVert g_t\rVert^2 .
\label{eq:deflection}
\end{equation}
For a target of constant amplitude this reduces to
$d^2 = T\,(A^2/\sigma^2)\,\lVert g\rVert^2$, so $d$ grows like $\sqrt{T}$: each
doubling of the number of frames moves the detection curve about $3$~dB toward
lower SNR. This is the familiar gain from integrating pulses over time
\cite{richards}, and the detection probability at a fixed false-alarm rate (FAR)
follows the Gaussian relation
\begin{equation}
P_d \;=\; Q\!\left(Q^{-1}(P_{\mathrm{FA}}) - d\right),
\label{eq:roc}
\end{equation}
where $Q$ is the standard normal tail. How far this gain extends is the central
theme of the thesis.

\paragraph{Correlated noise in a real sensor.} Real thermal cameras do not have
white noise. The standard three-dimensional (3-D) noise model \cite{dagostino}
separates sensor noise into parts that change from frame to frame and a fixed
pattern, the residual left after non-uniformity correction, that is the same in
every frame. The changing part averages out as $1/\sqrt{T}$; the fixed pattern
does not. When the noise is correlated, $n\sim\mathcal{N}(0,\Sigma)$, the optimal
detector becomes the whitening matched filter \cite{kay},
\begin{equation}
\Lambda(y) = s^\top \Sigma^{-1} y,
\qquad
d^2 = s^\top \Sigma^{-1} s,
\label{eq:whitening}
\end{equation}
which can still be computed when $\Sigma$ is known. Writing
$\Sigma = \sigma^2 I + C_{\mathrm{fix}}$, with $C_{\mathrm{fix}}$ the part that
repeats in every frame, leads to the claim the thesis sets out to prove:

\begin{quote}
For a slow or stationary target the signal lands on the same pixels in every
frame, where the fixed pattern also repeats, so the matching term in
\eqref{eq:whitening} does not vanish as $T$ grows and $d^2(T)$ stops increasing
at a finite floor. For a fast target the footprint moves to fresh pixels each
frame, the repeated part matters less, and the $\sqrt{T}$ growth lasts longer.
\end{quote}

\noindent Fixed-pattern noise therefore sets a limit on temporal integration
that depends on how fast the target moves. This is the limit Q3 asks us to map.

\section{Related Work}

Using the optimal detector to judge whether processing has kept a faint signal is
well established in medical imaging, where the ideal, or model, observer measures
task-based image quality \cite{barrett}. In other fields, deep networks are
compared with the matched filter at the output: gravitational-wave astronomy
\cite{gabbard}, radio-frequency signal detection with unknown parameters
\cite{anders}, and, recently, in-context hypothesis testing in transformers,
which links a network's internal representation to the likelihood-ratio statistic
\cite{chaudhry}. These studies either treat the network as a closed box, reading
only its output, or work in settings other than faint moving-target video.

A separate line of work looks inside networks. Linear classifier probes
\cite{alain} and information-plane analysis \cite{tishbyblackbox,ib} measure how
much task-relevant information each layer keeps. They do this for general tasks
and without an optimal-detector reference, so they cannot say how far a layer
sits from the best result physically possible.

On the temporal side, recent theory gives the longest span over which a recurrent
network can learn long-range structure \cite{livi}, a learning-based counterpart
to the classical $\sqrt{T}$ bound. Deep learning for infrared small-target
detection is itself a large and active area \cite{cheng}, but it competes on
accuracy and uses no optimal-detector reference.

What is missing is the combination of these pieces: the computable optimum used
inside the network, layer by layer, to find and repair where a faint signal is
lost, for moving-target infrared video, and tied to the temporal-integration
limit. The tool is borrowed from the ideal-observer tradition, but using it in
this way, and for this problem, appears to be new.

\section{Proposed Research}

\subsection{A testbed with a computable optimum}
We use a controllable synthetic infrared generator that implements the model in
\eqref{eq:model}: a Gaussian point target on a chosen trajectory, the 3-D
sensor-noise model \cite{dagostino} with its random and fixed-pattern parts, and
a multiplicative, time-correlated term for the target's brightness fluctuation.
Because the signal and the noise are known, the matched filter
\eqref{eq:deflection} and its whitening version \eqref{eq:whitening} give a
computable optimum $D^\star$ on the raw frames. The added realism is kept Gaussian
and of known covariance, so that this optimum stays exact (Section~\ref{sec:scope}).

\subsection{Finding where detectability is lost}
Write the detector as a composition $F = f_L\circ\cdots\circ f_1$, with
intermediate features $Z_k = (f_k\circ\cdots\circ f_1)(y)$. The label, the data,
and the representations form a Markov chain
$H \to y \to Z_1 \to \cdots \to Z_L \to \hat{H}$. Let $D_k$ be the best detection
performance, measured by deflection or by $P_d$ at a fixed FAR, that can be
obtained from $Z_k$ alone. The data-processing inequality then gives
\begin{equation}
I(Z_k;H) \le I(Z_{k-1};H) \;\Rightarrow\; D_k \le D_{k-1} \le \cdots \le D^\star,
\label{eq:dpi}
\end{equation}
so detectability can only stay level or fall as we move deeper, starting from the
optimum $D^\star$. Define the gap $G_k = D^\star - D_k$, the drop at stage $k$,
$\Delta D_k = D_{k-1} - D_k$, and the bottleneck stage
$k^\star = \arg\max_k \Delta D_k$, the one that loses the most. Each $D_k$ is
estimated with a fixed linear probe on the frozen features \cite{alain}. Since a
fixed probe can only read out part of what is present, this gives a lower bound,
and the size of the drop between stages is what we rely on. A plot of $D_k$
against depth, with $D^\star$ on top, is contribution C1, and the location of the
drop is C2, which we will trace as a function of SNR, sequence length, and motion.

\subsection{Repairing the bottleneck}
Once $k^\star$ is found, we replace $f_{k^\star}$ with a different
feature-reduction designed to keep the signal, leave the rest of the network
unchanged, and measure $D_k$ again. The repair has worked if the gap
$G_{k^\star}$ shrinks and detectability rises toward $D_{k^\star-1}$, and at the
output toward $D^\star$. The inequality in \eqref{eq:dpi} caps the possible gain:
no change to one stage can beat the optimum, so the claim can be tested directly
rather than only asserted. This is the precise and testable form of an earlier
idea, that a different way of collapsing the image might bring a lost target
back, now applied only at the stage where the loss is shown to occur.

\subsection{The temporal-integration limit}
Using \eqref{eq:whitening} and the claim in Section~2, we will trace $d^2(T)$ and
$P_d(T)$ for both the optimum and the learned detector, across SNR and target
speed, separating the white-noise $\sqrt{T}$ region from the fixed-pattern floor.
The network's limit is then placed against the learning-based horizon of
\cite{livi}: the classical bound says how much information $T$ frames hold, while
the learnable horizon says how much a trained recurrent network can actually use.
Comparing the two against the computable optimum is the synthesis the thesis aims
for.

\paragraph{Hypothesis.} We expect that at very low SNR the limiting stage is the
early spatial down-sampling, where a target only a few pixels wide is erased
before the network integrates over time, and that for longer sequences the
temporal stage becomes limiting instead. We also expect the fixed-pattern floor
to hurt slow targets most. Each part can be checked by the measurements above.

\section{Research Plan}

The work runs in stages, each depending on the one before it.
\begin{enumerate}
\item Extend the computable optimum from the white-noise matched filter to the
whitening matched filter \eqref{eq:whitening} for the 3-D-noise testbed.
\item Train a multi-frame detector (a per-frame U-Net encoder with a recurrent
temporal stage and detection heads) on the testbed, and record its detection
curve against $D^\star$.
\item Produce the detectability-against-depth plot at a low SNR where the optimum
still succeeds. A clear drop confirms an architecture-limited case and leads to
the repair step. A drop already at the input means the case is
information-limited, which is itself a useful result and tells us to shift the
SNR or sequence-length range.
\item Redesign the bottleneck stage and measure how much detectability returns,
bounded by \eqref{eq:dpi}.
\item Sweep the limit over SNR, sequence length, and motion, and relate it to the
learnable horizon.
\item Write up the results.
\end{enumerate}
The deliverables are the four contributions C1 to C4. Progress will be judged by
three figures: the detector against the optimum; the
detectability-against-depth plot showing the drop; and the before-and-after
comparison of the repair, together with the map of the temporal-integration limit.

\section{Scope and Limitations}
\label{sec:scope}
This is, by design, a study of performance limits on synthetic data. That choice
is what makes the optimum computable and the whole diagnosis possible, and it is
the main advantage over work on real footage. The price is a gap between the
synthetic data and reality. The noise parameters are set from physics rather than
fitted to a particular camera, since we do not assume access to real low-SNR
drone footage, and the added realism is kept Gaussian with known covariance so
that the optimum stays exact. Non-Gaussian clutter, sensor artifacts, and fitting
the noise model to real recordings are left for future work. Architectural
experiments that do not serve the main questions are also set aside.

\begin{thebibliography}{99}

\bibitem{marcum} J.~I.~Marcum, \emph{A Statistical Theory of Target Detection by
Pulsed Radar}, RAND Corporation Research Memoranda RM-753/RM-754, 1947--1948.

\bibitem{kay} S.~M.~Kay, \emph{Fundamentals of Statistical Signal Processing,
Vol.~II: Detection Theory}. Prentice-Hall, 1998.

\bibitem{dagostino} J.~A.~D'Agostino and C.~M.~Webb, ``Three-dimensional analysis
framework and measurement methodology for imaging system noise,'' \emph{Proc.
SPIE} 1488, \emph{Infrared Imaging Systems}, pp.~110--121, 1991.
doi:10.1117/12.45794.

\bibitem{barrett} H.~H.~Barrett, J.~Yao, J.~P.~Rolland, and K.~J.~Myers, ``Model
observers for assessment of image quality,'' \emph{Proc. Natl. Acad. Sci. USA},
vol.~90, no.~21, pp.~9758--9765, 1993. (See also H.~H.~Barrett and K.~J.~Myers,
\emph{Foundations of Image Science}, Wiley, 2004.)

\bibitem{gabbard} H.~Gabbard, M.~Williams, F.~Hayes, and C.~Messenger, ``Matching
matched filtering with deep networks for gravitational-wave astronomy,''
\emph{Phys. Rev. Lett.}, vol.~120, art.~141103, 2018. arXiv:1712.06041.

\bibitem{alain} G.~Alain and Y.~Bengio, ``Understanding intermediate layers using
linear classifier probes,'' arXiv:1610.01644, 2016.

\bibitem{tishbyblackbox} R.~Shwartz-Ziv and N.~Tishby, ``Opening the black box of
deep neural networks via information,'' arXiv:1703.00810, 2017.

\bibitem{ib} N.~Tishby, F.~C.~Pereira, and W.~Bialek, ``The information bottleneck
method,'' arXiv:physics/0004057, 1999.

\bibitem{livi} L.~Livi, ``Learnability window in gated recurrent neural
networks,'' arXiv:2512.05790, 2025.

\bibitem{anders} T.~Anders, H.~P.~Kothari, and R.~M.~Buehrer, ``On the ability of
deep learning to detect signals with unknown parameters,'' arXiv:2512.13542,
2025.

\bibitem{chaudhry} F.~Chaudhry and S.~Gadkari, ``Implicit statistical inference in
transformers: approximating likelihood-ratio tests in-context,''
arXiv:2603.10573, 2026.

\bibitem{cheng} Y.~Cheng, X.~Lai, Y.~Xia, and J.~Zhou, ``Infrared dim small target
detection networks: a review,'' \emph{Sensors}, vol.~24, no.~12, art.~3885, 2024.

\bibitem{richards} M.~A.~Richards, \emph{Fundamentals of Radar Signal
Processing}, 2nd ed. McGraw-Hill, 2014.

\end{thebibliography}

\end{document}
